\title{Nash-CaSH: Game-Theoretic Cross-Attention Harmonization for Training-Free Ultra-High-Resolution Video Generation}

\begin{abstract}
Generating ultra-high-resolution videos with pretrained video Diffusion Transformers is difficult because attention layers must operate far beyond the token scale observed during pretraining.
Local window attention improves detail fidelity by preserving a native-scale receptive field, yet it can weaken global semantic coordination and create repeated structures.
We propose Nash-CaSH, a training-free cross-attention harmonization strategy that treats local and global semantic evidence as two cooperative players in a Nash bargaining game.
At each cross-attention layer, Nash-CaSH computes branch payoffs from query-output compatibility and response confidence, then derives a closed-form Nash social-welfare allocation to fuse the global branch and local branch.
This adaptive equilibrium gate replaces fixed hand-tuned branch mixing and lets each token receive as much global semantic authority as it needs while retaining local detail when it is beneficial.
Nash-CaSH is compatible with cached global cross-attention, enabling efficient high-resolution inference without model finetuning.
\end{abstract}

\section{Introduction}

Video Diffusion Transformers have become a strong backbone for text-to-video generation, but their attention cost grows rapidly with spatial and temporal resolution.
A model pretrained at native resolution can often generate plausible low-resolution videos, yet direct high-resolution inference introduces a distribution shift in token count and receptive field.
The result is usually a trade-off: full attention preserves long-range structure but becomes expensive and may blur local details, while local attention improves texture and sharpness but can create repeated objects or inconsistent global layouts.

This work studies a training-free solution for high-resolution video synthesis.
The central idea is to preserve local native-scale evidence while still injecting global semantic coordination.
Instead of using a fixed mixing coefficient between the local and global branches, we formulate the interaction as a token-wise bargaining problem.
The local branch is responsible for training-scale detail fidelity; the global branch is responsible for semantic layout.
Their relative influence should vary across denoising steps, layers, prompts, and spatial positions.

We introduce Nash-CaSH, a Nash-equilibrium cross-attention harmonizer.
For every cross-attention output, Nash-CaSH estimates the utility of each branch and solves a Nash social-welfare allocation.
The resulting gate is adaptive, bounded, and closed-form, so it adds no trainable parameters and requires no optimization loop.
The method also preserves compatibility with feature caching: a cached global candidate can be reused across steps, while the current local candidate still participates in the equilibrium allocation.

\section{Method}

\subsection{Coarse-to-Fine High-Resolution Inference}

Given a prompt, the model first generates a base video at the native resolution of the pretrained backbone.
The video is then upsampled in pixel space, encoded into latents, and perturbed with Gaussian noise following an SDEdit-style refinement step.
High-resolution denoising begins from this noisy latent, preserving the base layout while allowing high-frequency details to be regenerated.

\subsection{Local-Global Dual Branches}

During high-resolution denoising, the latent is duplicated into two branches.
The local branch uses inward sliding-window self-attention, keeping the receptive field of each query close to the native training scale.
The global branch provides holistic context using full attention.
This design gives the model access to both local detail and global semantic structure.

Let the cross-attention outputs before output projection be
\begin{equation}
O_l(q), O_g(q) \in \mathbb{R}^{H \times D},
\end{equation}
where $q$ indexes a token, $l$ denotes the local branch, $g$ denotes the global branch, $H$ is the number of heads, and $D$ is the head dimension.

\subsection{Branch Payoffs}

Nash-CaSH measures the utility of each branch using two training-free signals: query compatibility and normalized response confidence.
For branch $b \in \{l,g\}$, with query $Q_b(q)$ and cross-attention output $O_b(q)$, the payoff is
\begin{equation}
u_b(q) =
\pi_b
\left[
\frac{1 + \cos(\mathrm{vec}(O_b(q)), \mathrm{vec}(Q_b(q)))}{2}
+
\sigma\left(
\frac{\mathrm{rms}(O_b(q))}{\mathrm{mean}_q\ \mathrm{rms}(O_b(q)) + \epsilon} - 1
\right)
+
\epsilon
\right].
\end{equation}
The cosine term rewards outputs aligned with the current query.
The response-confidence term rewards branch outputs that are strong relative to the layer's average response.
The prior $\pi_b$ can encode a conservative preference for global semantic stability.
The default values are $\pi_g=1.25$ and $\pi_l=1.0$.

\subsection{Nash Equilibrium Gate}

Let $a(q)$ be the semantic authority assigned to the global branch for token $q$; the local branch receives $1-a(q)$.
Nash-CaSH solves
\begin{equation}
a^*(q)
= \arg\max_{a \in [0,1]}
u_g(q) \log(a + \epsilon)
+ u_l(q) \log(1-a + \epsilon).
\end{equation}
This objective has the closed-form ratio
\begin{equation}
r(q)
=
\frac{u_g(q)^{1/\tau}}
{u_g(q)^{1/\tau} + u_l(q)^{1/\tau}},
\end{equation}
where $\tau$ is a temperature.
To avoid unstable extremes in training-free inference, the ratio is mapped into a bounded interval:
\begin{equation}
a(q) = a_{min} + (a_{max} - a_{min}) r(q).
\end{equation}
The fused cross-attention output is
\begin{equation}
O_{Nash}(q)
= a(q)O_g(q) + (1-a(q))O_l(q).
\end{equation}
The default interval is $a_{min}=0.65$, $a_{max}=0.98$.

\subsection{Cached Nash-CaSH}

Full global attention is expensive at high resolution.
Nash-CaSH supports cached global cross-attention by refreshing the global candidate every $P$ denoising steps.
For an intermediate step $t$, the fusion becomes
\begin{equation}
O_{Nash,t}(q)
= a_t(q)O_{g,t'}(q) + (1-a_t(q))O_{l,t}(q),
\end{equation}
where $t'$ is the most recent refresh step.
The local branch is still computed at the current step, so the gate can react to the current latent even when the global candidate is reused.

\section{Experiments}

The final submission should evaluate Nash-CaSH with fixed prompts, fixed seeds, official VBench metrics, and matched runtime conditions.
The table below is a reporting template; Nash-CaSH numbers should be filled only after official evaluation runs finish.

\begin{table}[t]
\centering
\caption{1080P evaluation template.}
\begin{tabular}{lcccccccc}
\hline
Method & Subj. & Back. & Motion & Aesth. & Image & Cons. & Score & Time \\
\hline
Native high-resolution inference & TBD & TBD & TBD & TBD & TBD & TBD & TBD & TBD \\
Fixed local-global override & TBD & TBD & TBD & TBD & TBD & TBD & TBD & TBD \\
Nash-CaSH w/o cache & TBD & TBD & TBD & TBD & TBD & TBD & TBD & TBD \\
Nash-CaSH w/ cache, P=2 & TBD & TBD & TBD & TBD & TBD & TBD & TBD & TBD \\
\hline
\end{tabular}
\end{table}

Required ablations include Nash gate disabled versus enabled, fixed override versus Nash-CaSH, bounds $a_{min}$ and $a_{max}$, branch priors, and cache interval $P$.

\section{Conclusion}

Nash-CaSH introduces a game-theoretic cross-attention harmonizer for training-free ultra-high-resolution video generation.
By converting local-global branch fusion into a Nash social-welfare allocation, the method adaptively balances semantic consistency and local detail without training or extra model parameters.
Its cache-compatible design keeps high-resolution inference practical.
Final numerical claims should be made only after completing the official evaluation protocol with fixed seeds and reproducible settings.
